\begin{theorem}[Remark~1.1(a): the scaling identity for every $t>0$]
Let $m\in\mathbb N$ and let $A:\mathbb R^{m}\to\mathbb R^{m\times m}$ and
$H:\mathbb R^{m}\to\mathbb R$ satisfy the structure conditions (E) and (S)
(\noderef{StructureConditionsES}). Write $\mathcal E_A$ for the ellipticity function of $A$
(\noderef{Ellip}). Then for every $t>0$ and every $\xi\in\mathbb R^{m}\setminus\{0\}$,
\[
 A(t\xi)=\frac{\mathcal E_A(t\xi)}{t^{2}\,\mathcal E_A(\xi)}\;A(\xi)
 \qquad\text{and}\qquad
 H(t\xi)=t\,\frac{\mathcal E_A(t\xi)}{t^{2}\,\mathcal E_A(\xi)}\;H(\xi).
\]
\end{theorem}
\begin{proof}
Let $\kappa:[1,\infty)\times(\mathbb R^{m}\setminus\{0\})\to(0,\infty)$ be a function for
which clause (6) of \noderef{StructureConditionsES} --- condition (S), (1.10) --- holds,
i.e.
\[
 A(t\xi)=\kappa(t,\xi)\,A(\xi)\quad\text{and}\quad H(t\xi)=t\,\kappa(t,\xi)\,H(\xi)
 \qquad\text{for all }t\ge1,\ \xi\ne0. \tag{S}
\]

\medskip\noindent\textbf{Step 0: two preliminary remarks.}

\emph{(0a) All denominators below are non-zero.} If $t>0$ and $\xi\ne0$ then $t\xi\ne0$,
because $t\xi=0$ with $t\ne0$ forces $\xi=0$. Hence condition (E) of
\noderef{StructureConditionsES}, which asserts $\mathcal E_A(\eta)>0$ for every $\eta\ne0$,
applies both at $\xi$ and at $t\xi$:
\[
 \mathcal E_A(\xi)>0,\qquad \mathcal E_A(t\xi)>0 . \tag{0a}
\]
In particular $t^{2}\mathcal E_A(\xi)>0$, so the quotient in the statement is a well-defined
positive real number.

\emph{(0b) The quadratic form is $2$-homogeneous in its vector slot.} Let
$M\in\mathbb R^{m\times m}$, $\eta\in\mathbb R^{m}$ and $s\in\mathbb R$. Since
$(s\eta)_k=s\,\eta_k$ for every $k$,
\[
 \sum_{i}\sum_{j}M_{i,j}\,(s\eta)_j\,(s\eta)_i
  =\sum_{i}\sum_{j}M_{i,j}\,s\,\eta_j\,s\,\eta_i
  =s^{2}\sum_{i}\sum_{j}M_{i,j}\,\eta_j\,\eta_i . \tag{0b}
\]
Applying this with $M=A(t\xi)$, $\eta=\xi$ and $s=t$, and recalling that
$\mathcal E_A(t\xi)$ is by \noderef{Ellip} the quadratic form of the matrix $A(t\xi)$
evaluated at the vector $t\xi$,
\[
 \mathcal E_A(t\xi)=\sum_{i}\sum_{j}A_{i,j}(t\xi)\,(t\xi)_j\,(t\xi)_i
  =t^{2}\sum_{i}\sum_{j}A_{i,j}(t\xi)\,\xi_j\,\xi_i . \tag{0b$'$}
\]

\medskip\noindent\textbf{Step 1: the case $t\ge1$, and the formula for $\kappa$.} Fix
$t\ge1$ and $\xi\ne0$. Substituting $A(t\xi)=\kappa(t,\xi)A(\xi)$ from (S) into
(0b$'$) and using that each entry satisfies
$A_{i,j}(t\xi)=\kappa(t,\xi)A_{i,j}(\xi)$,
\[
 \mathcal E_A(t\xi)=t^{2}\sum_{i}\sum_{j}\kappa(t,\xi)\,A_{i,j}(\xi)\,\xi_j\,\xi_i
  =t^{2}\,\kappa(t,\xi)\sum_{i}\sum_{j}A_{i,j}(\xi)\,\xi_j\,\xi_i
  =t^{2}\,\kappa(t,\xi)\,\mathcal E_A(\xi).
\]
By (0a) we may divide by $t^{2}\mathcal E_A(\xi)>0$:
\[
 \kappa(t,\xi)=\frac{\mathcal E_A(t\xi)}{t^{2}\,\mathcal E_A(\xi)} . \tag{1}
\]
Substituting (1) back into the two identities (S) gives, for every $t\ge1$ and $\xi\ne0$,
\[
 A(t\xi)=\frac{\mathcal E_A(t\xi)}{t^{2}\mathcal E_A(\xi)}A(\xi),
 \qquad
 H(t\xi)=t\,\frac{\mathcal E_A(t\xi)}{t^{2}\mathcal E_A(\xi)}H(\xi), \tag{2}
\]
which is the assertion of the theorem for $t\ge1$. (For $t=1$ formula (1) reads
$\kappa(1,\xi)=\mathcal E_A(\xi)/\mathcal E_A(\xi)=1$, so (2) is the trivial identity, as it
must be.)

\medskip\noindent\textbf{Step 2: the case $0<t<1$.} Fix $0<t<1$ and $\xi\ne0$, and set
\[
 \eta:=t\xi,\qquad t':=\frac1t .
\]
By Step 0(0a) applied to the pair $(t,\xi)$ we have $\eta\ne0$, and $t'>1$ because
$0<t<1$. So Step 1 is available for the pair $(t',\eta)$. Note first that
\[
 t'\eta=\frac1t\,(t\xi)=\xi,\qquad (t')^{2}=\frac1{t^{2}},\qquad
 \mathcal E_A(t'\eta)=\mathcal E_A(\xi),\qquad \mathcal E_A(\eta)=\mathcal E_A(t\xi),
\]
so that the scalar occurring in (2) for the pair $(t',\eta)$ is
\[
 \frac{\mathcal E_A(t'\eta)}{(t')^{2}\,\mathcal E_A(\eta)}
  =\frac{\mathcal E_A(\xi)}{t^{-2}\,\mathcal E_A(t\xi)}
  =\frac{t^{2}\,\mathcal E_A(\xi)}{\mathcal E_A(t\xi)} . \tag{3}
\]
By (0a) both $\mathcal E_A(\xi)>0$ and $\mathcal E_A(t\xi)>0$, so the scalar (3) is a
positive real number; call it $\lambda$. Then:

\emph{The identity for $A$.} The first identity of (2), applied to the pair
$(t',\eta)$ and read through $t'\eta=\xi$ and (3), says
\[
 A(\xi)=A(t'\eta)=\lambda\,A(\eta)=\lambda\,A(t\xi).
\]
Multiplying both sides by $\lambda^{-1}=\mathcal E_A(t\xi)/\big(t^{2}\mathcal E_A(\xi)\big)$,
which is legitimate because $\lambda>0$, gives
\[
 A(t\xi)=\lambda^{-1}A(\xi)=\frac{\mathcal E_A(t\xi)}{t^{2}\,\mathcal E_A(\xi)}\,A(\xi),
\]
as required.

\emph{The identity for $H$.} The second identity of (2), applied to the pair
$(t',\eta)$ and read through $t'\eta=\xi$, $t'=1/t$ and (3), says
\[
 H(\xi)=H(t'\eta)=t'\,\lambda\,H(\eta)
  =\frac1t\cdot\frac{t^{2}\mathcal E_A(\xi)}{\mathcal E_A(t\xi)}\cdot H(t\xi)
  =\frac{t\,\mathcal E_A(\xi)}{\mathcal E_A(t\xi)}\,H(t\xi).
\]
The scalar $t\,\mathcal E_A(\xi)/\mathcal E_A(t\xi)$ is positive by (0a) and $t>0$;
multiplying both sides by its reciprocal gives
\[
 H(t\xi)=\frac{\mathcal E_A(t\xi)}{t\,\mathcal E_A(\xi)}\,H(\xi)
  =t\cdot\frac{\mathcal E_A(t\xi)}{t^{2}\,\mathcal E_A(\xi)}\cdot H(\xi),
\]
as required.

\medskip\noindent\textbf{Step 3: conclusion.} Every $t>0$ satisfies either $t\ge1$, in
which case Step 1 gives both identities, or $0<t<1$, in which case Step 2 does. Since
$\xi\ne0$ was arbitrary, the two identities hold for all $t>0$ and all
$\xi\in\mathbb R^{m}\setminus\{0\}$.
\end{proof}
